% Generated by publication/build_sections.py; do not edit by hand.
\paragraph{Informal verdict.}
\textbf{A --- Complete solution}, with \textbf{high confidence}.

\paragraph{Lean coverage.}
Lean proves the full Q701 classification: for every function in the anchored Holder ball on [0,1], extremality is equivalent both to vanishing of the Holder-slack path pseudometric from 0 and to the finite-chain criterion with arbitrarily small total deficit. It also proves convexity and illustrative extreme and non-extreme examples. The result holds for arbitrary real alpha and therefore specializes to the printed range 0 < alpha < 1; optional strengthenings from the informal response are outside the canonical problem scope.

\subsection{Audited informal answer}
\subsubsection{1. Faithful restatement}

Fix an exponent \(\alpha\in(0,1)\). Let

\[
\mathcal C_\alpha
=
%
\left\{
f:[0,1]\to\mathbb R:
f(0)=0,\quad
|f(x)-f(y)|\le |x-y|^\alpha
\ \text{for all }x,y\in[0,1]
\right\}.
\]

Equipped with the uniform topology, \(\mathcal C_\alpha\) is a compact convex subset of \(C([0,1])\). The problem asks for simple properties and, ideally, a characterization of its extreme points. There are no subparts or additional qualifications in the PDF. 

The phrase ``\(\alpha\)-Hölder of ratio \(1\)'' must mean Hölder constant \textbf{at most} \(1\). If it meant that the optimal constant were exactly \(1\), the set would not be convex: \(x^\alpha\) and \(-x^\alpha\) would belong to it, whereas their midpoint \(0\) would not.

The answer below gives a complete intrinsic characterization.

\medskip\hrule\medskip

\subsubsection{2. The slack pseudometric}

Write

\[
d_\alpha(x,y)=|x-y|^\alpha.
\]

Since \(0<\alpha<1\), the inequality

\[
(r+s)^\alpha\le r^\alpha+s^\alpha
\]

shows that \(d_\alpha\) is a metric.

For \(f\in\mathcal C_\alpha\), define its \textbf{slack} on a pair ((x,y)) by

\[
\sigma_f(x,y)
=
%
= d_\alpha(x,y)-|f(x)-f(y)|
%
|x-y|^\alpha-|f(x)-f(y)|.
\]

Thus \(\sigma_f(x,y)\ge0\), and \(\sigma_f(x,y)=0\) precisely when the Hölder inequality is saturated at ((x,y)).

For a finite chain

\[
P=(x_0,x_1,\ldots,x_n)
\]

define its total slack by

\[
\ell_f(P)
=
%
\sum_{i=1}^{n}\sigma_f(x_{i-1},x_i).
\]

Finally, define

\[
\rho_f(x,y)
=
%
\inf
\left\{
\ell_f(x_0,\ldots,x_n):
n\ge1,\ x_0=x,\ x_n=y
\right\},
\]

and put \(\rho_f(x,x)=0\).

The function \(\rho_f\) is a pseudometric. Indeed, symmetry follows by reversing chains, and the triangle inequality follows by concatenating chains. Moreover, using the one-edge chain,

\[
0\le \rho_f(x,y)\le \sigma_f(x,y)\le d_\alpha(x,y).
\tag{2.1}
\]

It is also the largest pseudometric bounded above by \(\sigma_f\).

\medskip\hrule\medskip

\subsubsection{3. Complete characterization}

\paragraph{Theorem}

For \(f\in\mathcal C_\alpha\), the following conditions are equivalent.

\begin{enumerate}
\item \(f\) is an extreme point of \(\mathcal C_\alpha\).
\end{enumerate}

\begin{enumerate}
\item For every \(x\in[0,1]\),
\end{enumerate}

\[
   \rho_f(0,x)=0.
   \tag{3.1}
\]

\begin{enumerate}
\item For every \(x,y\in[0,1]\),
\end{enumerate}

\[
   \rho_f(x,y)=0.
   \tag{3.2}
\]

\begin{enumerate}
\item For every \(x\in[0,1]\) and every \(\varepsilon>0\), there exist an integer \(n\ge1\) and points
\end{enumerate}

\[
   0=x_0,x_1,\ldots,x_n=x
\]

   such that

\[
   \boxed{
   \sum_{i=1}^{n}
   \left(
   |x_i-x_{i-1}|^\alpha
   --------------------
%
   |f(x_i)-f(x_{i-1})|
   \right)
   <\varepsilon .
   }
   \tag{3.3}
\]

Equivalently, \(f\) is extreme precisely when every point can be joined to \(0\) by finite chains along which the accumulated deficit in the Hölder inequalities is arbitrarily small.

A general pointed-metric-space version of this criterion appears in a theorem of Jeff D. Farmer; the proof below is self-contained. ([SciSpace][1])

\medskip\hrule\medskip

\subsubsection{4. Proof of the characterization}

We first identify all possible midpoint perturbations of \(f\).

\paragraph{Lemma 4.1: midpoint perturbations}

Let \(u:[0,1]\to\mathbb R\) satisfy \(u(0)=0\). Then

\[
f+u\in\mathcal C_\alpha
\quad\text{and}\quad
f-u\in\mathcal C_\alpha
\]

if and only if

\[
|u(x)-u(y)|\le \sigma_f(x,y)
\qquad
\text{for all }x,y\in[0,1].
\tag{4.1}
\]

\subparagraph{Proof}

For real numbers (A,B),

\[
\max{|A+B|,\ |A-B|}=|A|+|B|.
\tag{4.2}
\]

Indeed, if \(A\) and \(B\) have the same sign, the first absolute value equals (|A|+|B|); if they have opposite signs, the second one does.

Apply this with

\[
A=f(x)-f(y),
\qquad
B=u(x)-u(y).
\]

The two conditions

\[
|A+B|\le d_\alpha(x,y),
\qquad
|A-B|\le d_\alpha(x,y)
\]

are therefore equivalent to

\[
|A|+|B|\le d_\alpha(x,y),
\]

which is exactly

\[
|u(x)-u(y)|
\le
d_\alpha(x,y)-|f(x)-f(y)|
=
%
\sigma_f(x,y).
\]

This proves the lemma. \(\square\)

\medskip\hrule\medskip

\paragraph{Lemma 4.2: path formulation}

Condition (4.1) is equivalent to

\[
|u(x)-u(y)|\le \rho_f(x,y)
\qquad
\text{for all }x,y.
\tag{4.3}
\]

\subparagraph{Proof}

If (4.1) holds, then for every chain \(x=x_0,\ldots,x_n=y\),

\[
\begin{aligned}
|u(x)-u(y)|
&\le
\sum_{i=1}^{n}|u(x_i)-u(x_{i-1})|\\
&\le
\sum_{i=1}^{n}\sigma_f(x_{i-1},x_i).
\end{aligned}
\]

Taking the infimum over all chains gives (4.3).

Conversely, (2.1) gives \(\rho_f(x,y)\le \sigma_f(x,y)\), so (4.3) implies (4.1). \(\square\)

Thus all symmetric decompositions

\[
f=\frac{(f+u)+(f-u)}2
\]

are described exactly by the \(1\)-Lipschitz functions for the pseudometric \(\rho_f\), anchored at \(0\).

\medskip\hrule\medskip

\paragraph{Proof that \(2\) implies extremality}

Suppose that \(\rho_f(0,x)=0\) for every \(x\).

Assume

\[
f=\frac{g+h}{2},
\qquad g,h\in\mathcal C_\alpha.
\]

Set

\[
u=\frac{g-h}{2}.
\]

Then \(g=f+u\), \(h=f-u\), and \(u(0)=0\). By Lemma 4.2,

\[
|u(x)|
=
%
|u(x)-u(0)|
\le
\rho_f(0,x)
=
%
0.
%
\]

Hence \(u=0\), so \(g=h=f\). Therefore \(f\) is extreme.

\medskip\hrule\medskip

\paragraph{Proof that failure of \(2\) implies non-extremality}

Suppose that

\[
\rho_f(0,x_*)>0
\]

for some \(x_*\in[0,1]\). Define

\[
u(x)=\rho_f(0,x).
\]

The reverse triangle inequality for the pseudometric \(\rho_f\) gives

\[
|u(x)-u(y)|
=
%
|\rho_f(0,x)-\rho_f(0,y)|
\le
\rho_f(x,y).
\]

By (2.1),

\[
|u(x)-u(y)|
\le \rho_f(x,y)
\le \sigma_f(x,y).
\]

Consequently, Lemma 4.1 gives

\[
f+u,\ f-u\in\mathcal C_\alpha.
\]

Moreover, \(u(0)=0\), while \(u(x_*)>0\), so \(u\neq0\). Thus

\[
f=\frac{(f+u)+(f-u)}2
\]

is a nontrivial midpoint decomposition, and \(f\) is not extreme.

This proves the equivalence of \(1\) and \(2\).

Conditions \(2\) and \(3\) are equivalent because

\[
\rho_f(x,y)
\le
\rho_f(x,0)+\rho_f(0,y).
\]

Finally, \(2\) and \(4\) are equivalent by the definition of the infimum \(\rho_f(0,x)\). The theorem is proved. \(\square\)

\medskip\hrule\medskip

\subsubsection{5. Quotient-space form of the answer}

Define an equivalence relation by

\[
x\sim_f y
\quad\Longleftrightarrow\quad
\rho_f(x,y)=0.
\]

The pseudometric \(\rho_f\) induces a genuine metric on the quotient

\[
X_f=[0,1]/!\sim_f.
\]

Every nontrivial decomposition of \(f\) is obtained by choosing a nonzero anchored \(1\)-Lipschitz function \(v\) on \(X_f\), pulling it back to \(u=v\circ q_f\), and writing

\[
f=\frac{(f+u)+(f-u)}2.
\]

Hence another exact formulation is

\[
\boxed{
f\text{ is extreme}
\quad\Longleftrightarrow\quad
X_f\text{ consists of a single point}.
}
\]

This not only determines whether \(f\) is extreme; it describes all its possible symmetric midpoint decompositions.

\medskip\hrule\medskip

\subsubsection{6. Saturated pairs and a simple sufficient condition}

Call (\{x,y\}) a \textbf{tight pair} if

\[
|f(x)-f(y)|=|x-y|^\alpha.
\]

Let \(G_f\) be the graph with vertex set ([0,1]) and an edge between every tight pair.

If every point can be connected to \(0\) by a finite path in \(G_f\), then every edge in the path has zero slack, so \(\rho_f(0,x)=0\). Thus \(f\) is extreme.

More generally, let \(R_f\) be the set of points connected to \(0\) by finite tight paths. If \(R_f\) is dense in ([0,1]), then \(f\) is extreme. Indeed, \(x\mapsto\rho_f(0,x)\) is \(\alpha\)-Hölder by

\[
|\rho_f(0,x)-\rho_f(0,y)|
\le \rho_f(x,y)
\le |x-y|^\alpha,
\]

and it vanishes on the dense set \(R_f\).

\paragraph{A stronger conclusion: exposed points}

If \(R_f\) is dense, then \(f\) is in fact an \textbf{exposed point} of \(\mathcal C_\alpha\).

To see this, choose a dense sequence \((z_n)\subset R_f\). For each \(n\), choose a finite tight path from \(0\) to \(z_n\). Enumerate all its oriented edges as \((p_m,q_m)\), orienting each one so that

\[
\varepsilon_m\bigl(f(q_m)-f(p_m)\bigr)
=
%
|q_m-p_m|^\alpha,
\qquad \varepsilon_m\in{-1,1}.
\]

Define the continuous linear functional

\[
\Lambda(g)
=
%
\sum_{m=1}^{\infty}
2^{-m}\varepsilon_m\bigl(g(q_m)-g(p_m)\bigr).
\]

For every \(g\in\mathcal C_\alpha\),

\[
\Lambda(g)
\le
\sum_{m=1}^{\infty}2^{-m}|q_m-p_m|^\alpha
=
%
\Lambda(f).
\]

Equality forces equality on every oriented tight edge, hence \(g(z_n)=f(z_n)\) for all \(n\). Continuity and density then give \(g=f\). Thus \(f\) is the unique maximizer of \(\Lambda\).

\medskip\hrule\medskip

\subsubsection{7. The finite-set classification}

Let

\[
F={t_0,t_1,\ldots,t_N},
\qquad t_0=0,
\]

be finite, and let \(\mathcal C_\alpha(F)\) be the analogous finite-dimensional polytope.

For \(f\in\mathcal C_\alpha(F)\), form the graph whose edges are precisely the saturated pairs

\[
|f(t_i)-f(t_j)|=|t_i-t_j|^\alpha.
\]

Then

\[
\boxed{
f\in\operatorname{Ext}\mathcal C_\alpha(F)
\quad\Longleftrightarrow\quad
\text{the saturated-pair graph is connected}.
}
\]

Indeed, in a finite graph the zero weighted-path distance from \(0\) to every vertex is equivalent to the existence of a path made entirely of zero-weight edges. If no such path exists, the minimum over the finitely many simple paths is strictly positive.

Equivalently, a finite extreme point is characterized by the existence of a spanning tree all of whose edge constraints are saturated.

This yields an exact \(O(N^2)\) recognition algorithm: construct all saturated edges and run a graph search from \(t_0\).

\medskip\hrule\medskip

\subsubsection{8. Explicit extreme points}

\paragraph{8.1 Distance cones}

For \(p\in[0,1]\), define

\[
k_p(x)=p^\alpha-|x-p|^\alpha.
\]

Then \(k_p(0)=0\). Since \(d_\alpha\) is a metric,

\[
|k_p(x)-k_p(y)|
=
%
\bigl|d_\alpha(p,x)-d_\alpha(p,y)\bigr|
\le d_\alpha(x,y),
\]

so \(k_p\in\mathcal C_\alpha\).

Moreover,

\[
k_p(p)-k_p(x)=|x-p|^\alpha.
\]

Thus every \(x\) is tightly connected to \(p\), and \(p\) is tightly connected to \(0\). Therefore \(k_p\) is extreme, indeed exposed. The same holds for \(-k_p\).

Special cases include

\[
k_0(x)=-x^\alpha,
\qquad
-k_0(x)=x^\alpha,
\]

and

\[
k_1(x)=1-(1-x)^\alpha.
\]

\medskip\hrule\medskip

\paragraph{8.2 A one-turn family}

For \(c\in[0,1]\), define

\[
T_c(x)=
\begin{cases}
x^\alpha, & 0\le x\le c,\\[2mm]
c^\alpha-(x-c)^\alpha, & c\le x\le1.
\end{cases}
\tag{8.1}
\]

We verify that \(T_c\in\mathcal C_\alpha\).

If both points lie on the same side of \(c\), this follows from

\[
|r^\alpha-s^\alpha|\le |r-s|^\alpha.
\]

Now let \(x\le c\le y\). Then

\[
\begin{aligned}
T_c(x)-T_c(y)
&=
x^\alpha-c^\alpha+(y-c)^\alpha\\
&\le (y-c)^\alpha
\le (y-x)^\alpha,
\end{aligned}
\]

while

\[
\begin{aligned}
T_c(y)-T_c(x)
&=
c^\alpha-x^\alpha-(y-c)^\alpha\\
&\le c^\alpha-x^\alpha
\le (c-x)^\alpha
\le (y-x)^\alpha.
\end{aligned}
\]

Hence \(|T_c(x)-T_c(y)|\le|x-y|^\alpha\).

Every \(x\le c\) is tightly connected to \(0\), since \(T_c(x)=x^\alpha\). Every \(x\ge c\) is tightly connected to \(c\), since

\[
T_c(c)-T_c(x)=(x-c)^\alpha.
\]

Thus \(T_c\) is extreme and exposed.

\medskip\hrule\medskip

\subsubsection{9. McShane-type construction of further extreme points}

Let \(A\subset[0,1]\) be compact, with \(0\in A\), and let \(g:A\to\mathbb R\) be an extreme anchored \(1\)-Lipschitz function for \(d_\alpha\).

Define the two canonical extensions

\[
M^+g(x)
=
%
\inf_{y\in A}\bigl(g(y)+|x-y|^\alpha\bigr),
\]

and

\[
M^-g(x)
=
%
\sup_{y\in A}\bigl(g(y)-|x-y|^\alpha\bigr).
\]

Both are elements of \(\mathcal C_\alpha\) and agree with \(g\) on \(A\).

Because \(A\) is compact, the infimum or supremum is attained. Thus for every \(x\), there is some \(y_x\in A\) such that either

\[
M^+g(x)-g(y_x)=|x-y_x|^\alpha
\]

or

\[
g(y_x)-M^-g(x)=|x-y_x|^\alpha.
\]

Consequently \(x\) is joined to \(A\) by a zero-slack edge. Within \(A\), extremality of \(g\) supplies arbitrarily small-slack chains to \(0\). Concatenating them proves, by the main theorem, that both \(M^+g\) and \(M^-g\) are extreme.

In particular, beginning with a finite extreme configuration---equivalently, a connected saturated graph---and applying either extension gives a large explicit class of extreme functions on the whole interval.

\medskip\hrule\medskip

\subsubsection{10. Necessary conditions, and why they are not sufficient}

Let

\[
[f]_\alpha
=
%
\sup_{x\ne y}
\frac{|f(x)-f(y)|}{|x-y|^\alpha}.
\]

\paragraph{Every extreme point lies on the Hölder sphere}

If \([f]_\alpha=L<1\), set

\[
u(x)=(1-L)x^\alpha.
\]

Then \([u]_\alpha=1-L\), and therefore

\[
[f\pm u]*\alpha
\le
[f]*\alpha+[u]_\alpha
=
%
1.
%
\]

Thus \(f\pm u\in\mathcal C_\alpha\), \(u\neq0\), and \(f\) is not extreme. Hence

\[
f\in\operatorname{Ext}\mathcal C_\alpha
\quad\Longrightarrow\quad
[f]_\alpha=1.
\tag{10.1}
\]

\paragraph{Hölder norm \(1\) is not sufficient}

Consider

\[
f(x)=x.
\]

Because \(0\le |x-y|\le1\),

\[
|x-y|\le |x-y|^\alpha,
\]

so \(f\in\mathcal C_\alpha\), and \([f]_\alpha=1\), with equality at ((0,1)).

Nevertheless \(f\) is not extreme. Define

\[
u(x)=(1-\alpha)x(1-x).
\]

For \(x<y\), put \(r=y-x\). Then

\[
|u(y)-u(x)|
=
%
(1-\alpha)r,|1-x-y|.
\]

Since \(0\le x\le1-r\),

\[
|1-x-y|\le1-r,
\]

and hence

\[
|u(y)-u(x)|
\le
(1-\alpha)r(1-r).
\]

The convexity inequality

\[
r^{-(1-\alpha)}
\ge
1+(1-\alpha)(1-r)
\]

gives, after multiplication by \(r\),

\[
r^\alpha-r
\ge
(1-\alpha)r(1-r).
\]

Therefore

\[
|u(y)-u(x)|
\le
r^\alpha-r
=
%
|x-y|^\alpha-|f(x)-f(y)|.
\]

Lemma 4.1 now yields

\[
f+u,\quad f-u\in\mathcal C_\alpha.
\]

Since \(u\neq0\), \(f\) is not extreme.

Thus the condition \([f]_\alpha=1\) is necessary but far from sufficient.

\medskip\hrule\medskip

\subsubsection{11. Strong extremality}

A point \(f\in\mathcal C_\alpha\) is called strongly extreme in the uniform norm if

\[
\left\lvert 
\frac{g_n+h_n}{2}-f
\right\rVert_\infty\longrightarrow0
\quad\Longrightarrow\quad
|g_n-h_n|*\infty\longrightarrow0
\]

for all sequences \(g_n,h_n\in\mathcal C_\alpha\).

Because \(\mathcal C_\alpha\) is compact in the uniform norm,

\[
\boxed{
f\text{ is extreme}
\quad\Longleftrightarrow\quad
f\text{ is strongly extreme}.
}
\]

Only the nontrivial implication needs proof. If it failed, there would be sequences with midpoint converging to \(f\) but

\[
|g_n-h_n|_\infty\ge\varepsilon>0.
\]

By compactness, a subsequence would satisfy

\[
g_n\to g,\qquad h_n\to h
\]

uniformly. Then

\[
f=\frac{g+h}{2},
\qquad
|g-h|_\infty\ge\varepsilon,
\]

contradicting extremality.

\medskip\hrule\medskip

\subsubsection{12. Boundary cases and necessity of the hypotheses}

\paragraph{The case \(\alpha=1\)}

The same slack-pseudometric theorem remains valid. It reduces to the classical concrete characterization

\[
\boxed{
f\in\operatorname{Ext}\mathcal C_1
\quad\Longleftrightarrow\quad
|f'(t)|=1
\ \text{for almost every }t\in[0,1].
}
\]

Here every \(1\)-Lipschitz function is absolutely continuous and \(|f'|\le1\) almost everywhere.

Indeed, for \(x<y\),

\[
\rho_f(x,y)
=
%
\int_x^y\bigl(1-|f'(t)|\bigr),dt.
\tag{12.1}
\]

For the upper bound, take increasingly fine monotone partitions. The sum of the absolute increments converges upward to the total variation

\[
\operatorname{Var}_{[x,y]}(f)
=
%
\int_x^y |f'(t)|,dt.
\]

For the reverse inequality, every edge \([z_{i-1},z_i]\) of an arbitrary chain satisfies

\[
|z_i-z_{i-1}|-|f(z_i)-f(z_{i-1})|
\ge
\int_{[z_{i-1},z_i]}
(1-|f'|).
\]

Every point between \(x\) and \(y\) is crossed by at least one edge of a chain from \(x\) to \(y\), so summation gives the reverse inequality in (12.1).

Notice the sharp change at the boundary: \(f(x)=x\) is extreme for \(\alpha=1\), but, as shown above, is not extreme for any \(0<\alpha<1\).

\paragraph{The case \(\alpha>1\)}

Every such Hölder function on an interval is constant. Indeed, subdividing ([x,y]) into \(n\) equal pieces gives

\[
|f(y)-f(x)|
\le
n\left(\frac{y-x}{n}\right)^\alpha
=
%
(y-x)^\alpha n^{1-\alpha}
\longrightarrow0.
\]

With \(f(0)=0\), the set consists only of the zero function, which is trivially extreme.

\paragraph{The case \(\alpha=0\)}

The stated compactness fails. Even among continuous functions, the family with oscillation at most \(1\) is not equicontinuous; for example,

\[
f_n(x)=\min(1,nx)
\]

has no uniformly convergent subsequence.

\paragraph{The anchoring condition}

The condition \(f(0)=0\) is essential. Without it, arbitrary constants may be added:

\[
f=\frac{(f+c)+(f-c)}2.
\]

Thus the unanchored set has no extreme points and is not bounded in the uniform norm.

\paragraph{The real-valued condition}

The scalar identity

\[
\max(|A+B|,|A-B|)=|A|+|B|
\]

is specifically real. For complex- or vector-valued Hölder maps, the midpoint perturbation condition and therefore the characterization require modification.

\medskip\hrule\medskip

\subsubsection{13. Compactness, separation, countability, and choice audit}

The ambient space is the real Banach space \(C([0,1])\) with the uniform norm, hence a Hausdorff metric space.

For every \(f\in\mathcal C_\alpha\),

\[
|f(x)|\le x^\alpha\le1,
\]

and the family has the common modulus of continuity \(r\mapsto r^\alpha\). It is closed under uniform convergence. The metric form of Arzelà--Ascoli therefore gives compactness. In this setting compactness and sequential compactness are equivalent.

The main extremality theorem itself is algebraic and metric; it uses neither compactness nor separability and works for arbitrary pointed metric spaces. Compactness is used only for the strong-extremality consequence and for attainment in the extension construction.

No general net argument, transfinite induction, or nonconstructive choice principle is required. The exposed-point argument uses only the separability of ([0,1]), through a countable dense sequence.

\medskip\hrule\medskip

\subsubsection{Final classification}

The extreme points of the compact convex set in Q701 are exactly the functions satisfying

\[
\boxed{
\forall x\in[0,1]\ \forall\varepsilon>0
\exists,0=x_0,x_1,\ldots,x_n=x
\quad
\sum_{i=1}^{n}
\left(
|x_i-x_{i-1}|^\alpha
--------------------
%
|f(x_i)-f(x_{i-1})|
\right)
<\varepsilon.
}
\]

Equivalently, the largest pseudometric lying below the pairwise Hölder slack

\[
|x-y|^\alpha-|f(x)-f(y)|
\]

must vanish identically. This criterion is necessary and sufficient and also gives a canonical nonzero perturbation \(u(x)=\rho_f(0,x)\) whenever \(f\) is not extreme.

[1]: https://scispace.com/pdf/extreme-points-of-the-unit-ball-of-the-space-of-lipschitz-3l772juamw.pdf "Extreme points of the unit ball of the space of Lipschitz functions"
